% Kapitel 6: Reflexion & Fazit
\section{Reflexion \& Fazit}

\subsection{Herausforderungen}

Die größte Herausforderung lag in der heterogenen Struktur der Zielwebsite. Dynamisch geladene Inhalte und unterschiedliche Seitenformate erschwerten das konsistente Extrahieren von Textinhalten. Der Web-Crawler musste mehrfach angepasst werden.

Eine weitere Herausforderung war die Chunk-Strategie: Die initiale rein zeichenbasierte Aufteilung führte zu Chunks, die mitten im Satz endeten. Durch Satzgrenzen-Erkennung konnte die Qualität verbessert werden.

In der UI-Phase mussten ergonomische Anforderungen mit visueller Klarheit abgewogen werden. Die Session-Verwaltung durfte nicht zu leeren Chats führen (lazy creation), und die Spracheingabe musste ohne zusätzliche Server-Infrastruktur nutzbar bleiben.

\subsection{Limitationen}

\begin{itemize}
    \item \textbf{Stichprobengröße:} N=10 ist für formative Tests ausreichend, erlaubt aber nur begrenzte Generalisierbarkeit.
    \item \textbf{Zeitliche Diskrepanz:} Vorschlagsfragen, Sessions, Spracheingabe und Theme wurden nach der Evaluation ergänzt und sind nicht durch SUS/Amershi abgesichert.
    \item \textbf{Lokales LLM:} Ollama mit Qwen 2.5 liefert gute Ergebnisse, ist aber langsamer als Cloud-Modelle und qualitativ weniger stabil bei komplexen Anfragen.
    \item \textbf{Einfaches Chunking:} Die zeichenbasierte Aufteilung ist nicht tokenbewusst.
    \item \textbf{Keine automatische Nachlern-Funktion:} Neue Informationen erfordern einen vollständigen Re-Import.
    \item \textbf{Browserabhängige Spracheingabe:} Die Web Speech API ist nicht in allen Browsern gleich verfügbar.
\end{itemize}

\subsection{Verbesserungsvorschläge}

Folgeevaluation der erweiterten UI; Synonym-Erkennung für Fachbegriffe; tokenbasiertes Chunking; WCAG-Prüfung.

\subsection{Fazit}

Ergonomia zeigt, dass ein RAG-Chatbot mit lokalen Komponenten nutzerzentrierte Antworten auf Basis einer Website liefern kann (SUS 73,5). Human-Centred Design nach ISO 9241-210~\cite{iso9241-210} prägte den gesamten Prozess von Interviews über Architekturentscheidungen bis zur Evaluation. Ergonomia ist ein UX-Projekt, das einen Chatbot als Interaktionsmedium nutzt.
